\section{The three investment vehicles}\label{sec:models}

Throughout the paper we model a residential building made of $n$
independent rental units, observed over a horizon $T$ years. The building's
total par value at inception is $P_0$. Each apartment $i \in \{1, \ldots,
n\}$ has a stochastic default time $\tau_i$ with cumulative probability
$\Prob(\tau_i \le T) = p$ matched to a calibrated marginal hazard. The
fraction $\ell_i \in [0, 1]$ of par lost when apartment $i$ defaults is
drawn from a Beta distribution. The aggregate cumulative loss path is
\begin{equation}
L(t) = \frac{1}{n}\sum_{i = 1}^{n} \ell_i \,\mathbb{1}\{\tau_i \le t\},
\qquad t \in [0, T].
\label{eq:agg-loss}
\end{equation}
The building generates a net rental flow $R_t = (g - m)\, S_t\,(1 - L(t))$
in monetary units, where $S_t$ is the market value, $g$ is the gross
rental yield and $m$ the maintenance cost share, both expressed per year.

\paragraph{Model A — Classical full ownership.}
A single investor owns one apartment of the building (par $P_0 / n$). Cash
flows in period $[t, t + \Delta t]$ are $R_t\, \Delta t / n$ while
$\tau_1 > t$ and $0$ thereafter. At $T$ the apartment is sold at the
prevailing market value $S_T / n$ regardless of any past default.

\paragraph{Model B — SAS shareholders.}
$n$ investors each buy one share of a SAS that holds the whole building.
Each share receives the pro-rata net rent
$R_t\, \Delta t / n$ each period and $S_T (1 - L(T)) / n$ at terminal.

\paragraph{Model C — CDO-like tranches.}
The same aggregate cash flow $R_t$ is partitioned across a stack of
tranches $\tau = (\tau_\text{equity}, \tau_\text{mezz}, \tau_\text{senior})$
with attach / detach points $(0, 0.25)$, $(0.25, 0.60)$ and $(0.60, 1.00)$.
Tranche $\tau$ absorbs the cumulative loss via the stop-loss payoff
\begin{equation}
L_\tau(t) = \min\bigl(\max(L(t) - a_\tau, 0), d_\tau - a_\tau\bigr),
\label{eq:stop-loss}
\end{equation}
and the periodic net rent is allocated waterfall-style: senior coupon
first, then mezzanine, then equity collects the residual. At terminal the
sale proceeds $S_T (1 - L(T))$ are paid back to surviving notionals in the
same order.

\begin{assumption}\label{ass:par}
Total portfolio par is $P_0 = \nObligors \times 45\,\text{m}^2 \times
10\,000\,\EUR / \text{m}^2 = \parTotal$ for the Paris-intermediate scenario
that drives every numerical result reported below.
\end{assumption}
